\documentclass[11pt,a4paper]{article}
\usepackage[utf8]{inputenc}
\usepackage[T1]{fontenc}
\usepackage[margin=0.75in]{geometry}
\usepackage{enumitem}
\usepackage{titlesec}
\usepackage{xcolor}
\usepackage{hyperref}

\definecolor{primarycolor}{RGB}{33,37,41}
\definecolor{secondarycolor}{RGB}{0,102,204}
\definecolor{accentcolor}{RGB}{108,117,125}
\hypersetup{colorlinks=true,linkcolor=secondarycolor,urlcolor=secondarycolor,citecolor=secondarycolor}
\titleformat{\section}{\Large\bfseries\color{primarycolor}}{}{0em}{}
\titlespacing*{\section}{0pt}{1.1em}{0.45em}
\setlist[itemize]{leftmargin=*,itemsep=0.2em,topsep=0.25em}
\pagestyle{empty}

\begin{document}

\begin{center}
{\Huge\bfseries\color{primarycolor} Marta Kowal-Brandt}\\[0.35em]
{\large\color{accentcolor} DevOps / Platform Engineering | QA Automation & Test Infrastructure}
\end{center}
\vspace{0.4em}

\section*{PROFESSIONAL SUMMARY}
QA automation engineer with experience owning GitLab CI pipelines and ephemeral environments for three teams at Odrava Health Software. Built Kubernetes-based merge-request environments and improved CI pipelines, reducing average runtime from 38 to 14 minutes and monthly CI runner costs by about 30\%. Brings hands-on experience with Helm, Docker, Kubernetes, GitLab CI, and observability, alongside personal k3s/GitOps and Terraform projects. Currently transitioning toward DevOps and platform engineering.

\section*{SKILLS}
CI/CD: GitLab CI, Jenkins, pipeline optimization, automated regression and end-to-end testing\\
Containers and orchestration: Docker, Kubernetes (namespace-level access; Helm charts for test environments), Helm, k3s, Argo CD\\
Infrastructure and automation: Terraform (personal Hetzner VPS project), Ansible (basic), Bash\\
Observability and data: Prometheus, Grafana, Loki, PostgreSQL\\
Programming and testing: Python, Java, pytest, pytest-xdist, Selenium, GitLab API\\
Cloud learning: AWS EC2 and S3 basics (personal free-tier use)

\section*{EXPERIENCE}
\textbf{QA Automation Engineer (Test Infrastructure) | Odrava Health Software | 2021 to Present}
\begin{itemize}
\item Own GitLab CI pipelines and ephemeral environments for three teams.
\item Built merge-request environments with an application, PostgreSQL, and seeded data, deployed via Helm to a shared Kubernetes cluster from GitLab CI; environments are removed after merge or 48 hours of inactivity.
\item Profiled pipelines across three teams, parallelized test jobs, added dependency caching, and split end-to-end tests into a separate stage; reduced average pipeline time from 38 to 14 minutes and monthly CI runner costs by about 30\%.
\item Built a flaky-test dashboard using CI results, PostgreSQL, Python, Grafana, and the GitLab API; the flaky-test rate fell from 7\% to below 1\% after teams addressed the leading offenders.
\item Wrote the internal guide to how CI works, which is used in onboarding.
\item Mentored two QA interns in 2023.
\item Joined two incident calls as an observer.
\end{itemize}

\textbf{QA Engineer | Brightmill Software House | 2019 to 2021}
\begin{itemize}
\item Tested web and mobile applications for different clients.
\item Introduced the first automated regression suite on the largest project, using Java and Selenium, and configured Jenkins to run it nightly.
\item The full regression had previously taken two people three days.
\end{itemize}

\section*{SELECTED PROJECT}
\textbf{Homelab k3s Cluster with GitOps | Personal | 2024 to 2025}
\begin{itemize}
\item Built a three-node Raspberry Pi 5 k3s cluster with applications deployed from a Git repository using Argo CD.
\item Set up Prometheus, Grafana, and Loki for monitoring, Traefik ingress, and cert-manager for TLS.
\item Provisioned infrastructure for a small Hetzner VPS with Terraform.
\item Technologies: k3s, Argo CD, Prometheus, Grafana, Loki, Terraform, Ansible (basic).
\end{itemize}

\section*{CERTIFICATIONS AND DEVELOPMENT}
ISTQB Foundation Level, 2019\\
Certified Kubernetes Administrator (CKA): studying; exam booked for January 2027.\\
Terraform Associate: completed practice tests; exam not taken.

\end{document}